\naslov{Krčenje razsežnosti}

Prvo vprašanje je, zakaj bi razsežnost sploh krčili.
Odgovor je prekletstvo večrazsežnosti:
Naj bodo $X_1, \ldots, X_n$ točke v enotski krogli v $\R^p$, porazdeljene
enakomerno, in naj bo $M$ najmanjša razdalja teh točk do središča.
Izračunamo lahko
\[
  P(M \le x) = 1 - (1 - x^p)^n.
\]
Potem je mediana (tj.~razdalja, pri kateri je najbližja točka z verjetnostjo
$\pol$ v krogli s tem radijem) enaka
\[
  x_m = F_M^{-1}(\pol) = \left( 1 - \inv{2^{1/n}} \right)^{1/p}.
\]
Z veliko točkami se oddaljujemo od $1$, z veliko dimenzijami pa se bližamo $1$.
Pozorni moramo biti na razmerje $n/p$.

\vprasanje{Demonstriraj problem večrazsežnosti.}

Želeli bi si \pojem{rangiranje} napovednih spremenljivk $X_i$ po pomembnosti.
Za to definiramo \pojem{relevantnost} spremenljivke kot funkcijo $r_S : \mb{X}
\to \R$, ki za podatkovno množico $S$ pove, kako pomembna je spremenljivka $X_i
\in \mb{X}$ (tu je $\mb{X}$ množica spremenljivk).

Mere relevantnosti ločimo na nadzorovane in nenadzorovane.
Za nenadzorovane opazimo, da je relevantnost je pozitivno korelirana z varianco
spremenljivke, saj večja varianca pomeni, da imamo več prostora za odločitve.
Ta mera pa se izgubi, če podatke normaliziramo.
V primeru diskretne spremenljivke lahko namesto variance vzamemo entropijo
\[
  H(X_i) = - \sum_{v} p_v \log_2 p_v,
\]
kjer je $p_v = \abs{\{x_i \in S \such x_i = v\}} / \abs{S}$.

\vprasanje{Opiši nenadzorovano mero relevantnosti.}

Pri nadzorovanih merah relevantnosti začnemo z uni-variantnimi metodami.
Pri teh gledamo samo povezanost med $X_i$ in $Y$ za nek $i$.
Ker uporabljajo $Y$, so nadzorovane; ker ignorirajo vse ostale $X_j$, so
uni-variantne.
Če sta $X_i$ in $Y$ numerični, je ena možnost korelacijski faktor, imamo pa še
cel kup drugih možnosti.
V primeru numerične $X_i$ in diskretne $Y$ lahko opazujemo \pojem{relief}, kjer
za naključno izbran primer $(\mb{x},y) \in S$ poiščemo najbližjega soseda v
projekciji na $X_i$, $\mb{x}_H$, ki pripada istemu razredu kot $\mb{x}$, in
najbližjega soseda $\mb{x}_M$, ki pripada različnemu razredu.
Če je razdalja do $\mb{x}_M$ veliko večja od razdalje do $\mb{x}_H$, je to dober
znak za relevantnost spremenljivke, na katero smo projecirali.
V algoritmu je relevantnost $X_i$ potem enaka
\[
  \inv{\abs{S}} \sum_{(\mb{x}, y) \in S} (\abs{\mb{x}_{Mi} - \mb{x}_i} -
  \abs{\mb{x}_{Hi} - \mb{x}_i}).
\]

\vprasanje{Opiši metodo reliefa.}

Zadnja možnost je ocenjevanje relevantnosti na podlagi modela.
Če imamo razumljiv model, npr.~linearni, lahko relevantnost direktno preberemo
iz njega.
V primeru manj razumljivega modela lahko permutiramo vrednosti v stolpcu $i$ v
podatkih in natreniramo nov model; spremenljivka je tako uničena, relevantnost
dobimo kot razliko v napovedni napaki.
Ta metoda je najpočasnejša, a je zelo kvalitetna.

\vprasanje{Kako oceniš relevantnost spremenljivke z danim modelom?}

Pri zmanjševanju razsežnosti pogosto želimo tvoriti nove spremenljivke, ki bodo
bolj relevantne za nas.
Običajno ima to obliko linearne transformacije $Z = X W$, kjer je $Z$ dimenzije
$q \ll p$.
Za to imamo metodo glavnih komponent, kjer dodatno zahtevamo, da je $W$
ortogonalna, in da dimenzije $Z$ pojasnijo čim več variance v osnovnih podatkih.
Predpostavimo, da imajo vse spremenljivke v $X$ povprečje $0$, torej je
kovariančna matrika enaka $C = X^T X$.
Ta je simetrična nenegativno definitna, torej jo lahko s pomočjo singularnega
razcepa zapišemo kot $C = W D W^T$, kjer je $D$ diagonalna matrika kvadratov
singularnih vrednosti.

\vprasanje{Razloži metodo glavnih komponent.}

% LocalWords:  rangiranje korelirana uni-variantnimi uni-variantne projecirali
